% Appendix B, printed pp. 73--80. Basic intersection theory used in Section 6.
\section{Basic intersection theory}\label{kps-sec-intersection}

\begin{definition}[The Grothendieck group of proper-support sheaves]\label{kps-def-K}
\dcref{B.1}\source{kleiman-picard:page-0073}
Let $\mathcal F(X/S)$ be the Abelian category of coherent sheaves whose
support is proper over a zero-dimensional Noetherian (Artin) subscheme of
$S$, and let $\mathcal F_r$ be the subcategory with support of dimension at
most $r$. Let $\mathcal K(X/S)$ be the free Abelian group on these sheaves
modulo short exact sequences, with filtration $\mathcal K_r$ generated by
$\mathcal F_r$. Euler characteristic induces $\chi:\mathcal K\to\mathbb Z$.
For an invertible $\mathcal L$, define
\[
 c_1(\mathcal L)\mathcal F=\mathcal F-\mathcal L^{-1}\ox\mathcal F.
\]
\end{definition}

\begin{lemma}[Divisor action]\label{kps-lem-divisor-action}
\dcref{B.2}\source{kleiman-picard:page-0073}
If $Y\subset X$ has $\mathcal O_Y$ in $\mathcal F$ and $D\subset Y$ is an
effective divisor with $\mathcal O_Y(D)\simeq\mathcal L|_Y$, then
$c_1(\mathcal L)[Y]=[D]$.
\uses{kps-def-K}
\end{lemma}
\begin{proof}
The exact sequence $0\to\mathcal O_Y(-D)\to\mathcal O_Y\to\mathcal O_D\to0$
identifies the right side with
$\mathcal O_Y-\mathcal L^{-1}\ox\mathcal O_Y=c_1(\mathcal L)[Y]$.
\end{proof}

\begin{lemma}[Additivity of first Chern operators]\label{kps-lem-c1-additivity}
\dcref{B.3}\source{kleiman-picard:page-0073}
For invertible $\mathcal L,\mathcal M$,
\[
c_1(\mathcal L)c_1(\mathcal M)=c_1(\mathcal L)+c_1(\mathcal M)-c_1(\mathcal L\ox\mathcal M),
\quad c_1(\mathcal L)c_1(\mathcal L^{-1})=c_1(\mathcal L)+c_1(\mathcal L^{-1}),
\quad c_1(\mathcal O_X)=0.
\]
The operators commute.
\uses{kps-def-K}
\end{lemma}
\begin{proof}
Apply both sides to a class $\mathcal F$ and expand; every side is
$\mathcal F-\mathcal L^{-1}\ox\mathcal F-\mathcal M^{-1}\ox\mathcal F+
\mathcal L^{-1}\ox\mathcal M^{-1}\ox\mathcal F$. The other assertions follow
from this identity and commutativity of tensor products.
\end{proof}

\begin{lemma}[Cycle decomposition]\label{kps-lem-cycle-decomposition}
\dcref{B.4}\source{kleiman-picard:page-0074}
For $\mathcal F\in\mathcal F_r$, let $Y_i$ be the $r$-dimensional
irreducible components of its support with reduced structure, and let $l_i$
be the generic stalk lengths. Then
\[
 \mathcal F\equiv\sum_i l_i[Y_i]\pmod{\mathcal K_{r-1}}.
\]
\uses{kps-def-K}
\end{lemma}
\begin{proof}
After restricting to a Noetherian neighborhood, the class of sheaves for
which the formula holds is closed under short exact sequences because generic
length is additive. It contains the structure sheaf of every integral closed
subscheme, so dévissage gives all of $\mathcal F_r$.
\end{proof}

\begin{lemma}[Dimension drop]\label{kps-lem-dimension-drop}
\dcref{B.5}\source{kleiman-picard:page-0074}
For every $r$, $c_1(\mathcal L)\mathcal K_r\subset\mathcal K_{r-1}$.
\uses{kps-lem-cycle-decomposition}
\end{lemma}
\begin{proof}
Tensoring by $\mathcal L^{-1}$ does not change the generic stalk at any
$r$-dimensional support component. The cycle decomposition lemma therefore
cancels the top-dimensional terms in
$\mathcal F-\mathcal L^{-1}\ox\mathcal F$.
\end{proof}

\begin{lemma}[Binomial expansion]\label{kps-lem-binomial-expansion}
\dcref{B.6}\source{kleiman-picard:page-0074}
For $\mathcal F\in\mathcal K_r$ and $m\in\mathbb Z$,
\[
 \mathcal L^m\ox\mathcal F=
 \sum_{i=0}^r\binom{m+i-1}{i}c_1(\mathcal L)^i\mathcal F.
\]
\uses{kps-lem-dimension-drop}
\end{lemma}
\begin{proof}
The formal identity
$y^m=\sum_{i\ge0}\binom{m+i-1}{i}(1-y^{-1})^i$ becomes the displayed
identity after setting $y=\mathcal L$. Terms with $i>r$ vanish by dimension
drop.
\end{proof}

\begin{theorem}[Snapper polynomial]\label{kps-thm-snapper}
\dcref{B.7}\source{kleiman-picard:page-0074}
For $\mathcal L_1,\ldots,\mathcal L_n\in\Pic(X)$, integers $m_i$, and
$\mathcal F\in\mathcal K_r$, the Euler characteristic
$\chi(\mathcal L_1^{m_1}\ox\cdots\ox\mathcal L_n^{m_n}\ox\mathcal F)$
is a polynomial in the $m_i$ of total degree at most $r$, explicitly
\[
 \sum_{i_1+\cdots+i_n\le r}\chi(c_1(\mathcal L_1)^{i_1}\cdots
 c_1(\mathcal L_n)^{i_n}\mathcal F)
 \prod_j\binom{m_j+i_j-1}{i_j}.
\]
\uses{kps-lem-binomial-expansion,kps-lem-dimension-drop}
\end{theorem}
\begin{proof}
Apply the binomial expansion successively to each tensor factor and use the
dimension-drop lemma to truncate the resulting sum.
\end{proof}

\begin{definition}[Intersection numbers]\label{kps-def-intersection-number}
\dcref{B.8}\source{kleiman-picard:page-0075}
For $\mathcal L_1,\ldots,\mathcal L_r$ and $\mathcal F\in\mathcal K_r$ set
\[
 \int c_1(\mathcal L_1)\cdots c_1(\mathcal L_r)\mathcal F
 :=\chi(c_1(\mathcal L_1)\cdots c_1(\mathcal L_r)\mathcal F).
\]
When $\mathcal F=\mathcal O_X$, omit it; for divisor line bundles use the
usual $(D_1\cdots D_r\cdot\mathcal F)$ notation.
\uses{kps-def-K}
\end{definition}

\begin{theorem}[First intersection properties]\label{kps-thm-first-intersection}
\dcref{B.9}\source{kleiman-picard:page-0075}
The intersection symbol vanishes if $\mathcal F\in\mathcal F_{r-1}$, is
symmetric and additive in each line bundle and in $\mathcal F$, and satisfies
\[
 \int c_1(\mathcal L_1)\cdots c_1(\mathcal L_r)\mathcal F
 =\sum_{i=0}^r(-1)^i\chi((\bigwedge^i\mathcal E)\ox\mathcal F),
 \quad\mathcal E=\bigoplus_j\mathcal L_j^{-1}.
\]
\uses{kps-lem-c1-additivity,kps-lem-dimension-drop,kps-def-intersection-number}
\end{theorem}
\begin{proof}
Dimension drop gives the vanishing. Additivity and symmetry follow from the
operator identities in Lemma B.3, and the exterior-power formula is the
expanded definition of the product of the commuting operators.
\end{proof}

\begin{corollary}[Surface intersection expansion]\label{kps-cor-surface-expansion}
\dcref{B.10}\source{kleiman-picard:page-0075}
For $\mathcal F\in\mathcal K_2$,
\[
\int c_1(\mathcal L_1)c_1(\mathcal L_2)\mathcal F
 =\chi(\mathcal F)-\chi(\mathcal L_1^{-1}\ox\mathcal F)-\chi(\mathcal L_2^{-1}\ox\mathcal F)
 +\chi(\mathcal L_1^{-1}\ox\mathcal L_2^{-1}\ox\mathcal F).
\]
\uses{kps-thm-first-intersection}
\end{corollary}
\begin{proof}
This is the $r=2$ case of the exterior-power formula in Theorem B.9.
\end{proof}

\begin{corollary}[Cohen--Macaulay intersection formula]\label{kps-cor-CM-intersection}
\dcref{B.11}\source{kleiman-picard:page-0075}
If effective divisors $D_1,\ldots,D_r$ meet the support of a coherent
$\mathcal F\in\mathcal F_r$ in finitely many points where $\mathcal F$ is
Cohen--Macaulay, then
\[
 (D_1\cdots D_r\cdot\mathcal F)=\length H^0(\mathcal F\ox\mathcal O_{D_1\cap\cdots\cap D_r}).
\]
\uses{kps-thm-first-intersection}
\end{corollary}
\begin{proof}
Use the Koszul complex of the defining sections. Cohen--Macaulayness makes
the sequence regular on the support, so higher Koszul homology vanishes. The
exterior-power formula and additivity of Euler characteristic leave the
length of the degree-zero homology.
\end{proof}

\begin{lemma}[Intersection cycles]\label{kps-lem-intersection-cycles}
\dcref{B.12}\source{kleiman-picard:page-0075}
If $\mathcal F\in\mathcal F_r$ has $r$-dimensional support components $Y_i$
with generic lengths $l_i$, then
\[
 \int c_1(\mathcal L_1)\cdots c_1(\mathcal L_r)\mathcal F
 =\sum_i l_i\int c_1(\mathcal L_1)\cdots c_1(\mathcal L_r)[Y_i].
\]
\uses{kps-lem-cycle-decomposition,kps-thm-first-intersection}
\end{lemma}
\begin{proof}
Apply the cycle decomposition lemma and the vanishing and additivity parts of
Theorem B.9.
\end{proof}

\begin{lemma}[Intersection with a divisor]\label{kps-lem-intersection-divisor}
\dcref{B.13}\source{kleiman-picard:page-0075}
If $D\subset Y$ is an effective divisor with $\mathcal O_Y(D)\simeq
\mathcal L_r|_Y$, then
\[
 \int c_1(\mathcal L_1)\cdots c_1(\mathcal L_r)[Y]
 =\int c_1(\mathcal L_1)\cdots c_1(\mathcal L_{r-1})[D].
\]
\uses{kps-lem-divisor-action}
\end{lemma}
\begin{proof}
Apply the divisor-action lemma to the last factor and then use the definition
of the intersection symbol.
\end{proof}

\begin{proposition}[Positivity for ample bundles]\label{kps-prop-ample-positivity}
\dcref{B.14}\source{kleiman-picard:page-0076}
If all $\mathcal L_j$ are ample and $\mathcal F\in\mathcal F_r$ is not in
$\mathcal K_{r-1}$, then
\[
 \int c_1(\mathcal L_1)\cdots c_1(\mathcal L_r)\mathcal F>0.
\]
\uses{kps-lem-intersection-cycles,kps-lem-intersection-divisor,kps-thm-first-intersection}
\end{proposition}
\begin{proof}
Reduce by intersection cycles to $\mathcal F=\mathcal O_Y$ for an integral
$r$-fold. Replace the last ample bundle by a very ample power and intersect
with a nonempty hyperplane divisor. The divisor-intersection lemma and
induction on $r$ give positivity; the case $r=0$ is the positive length of
nonzero $H^0(\mathcal F)$.
\end{proof}

\begin{lemma}[Functorial intersection pushforward]\label{kps-lem-intersection-pushforward}
\dcref{B.15}\source{kleiman-picard:page-0076}
For an $S$-map $g:X'\to X$, line bundles $\mathcal L_j$ on $X$, and
$\mathcal F\in\mathcal F_r(X'/S)$,
\[
 \int c_1(g^*\mathcal L_1)\cdots c_1(g^*\mathcal L_r)\mathcal F
 =\int c_1(\mathcal L_1)\cdots c_1(\mathcal L_r)\,g_*\mathcal F.
\]
\uses{kps-def-K,kps-thm-first-intersection}
\end{lemma}
\begin{proof}
Derived pushforward preserves the filtered Grothendieck groups and Euler
characteristic, and the projection formula commutes with each $c_1$. Higher
direct images have support of dimension at most $r-1$, so their classes do not
affect the $r$-fold intersection by Theorem B.9.
\end{proof}

\begin{proposition}[Projection formula]\label{kps-prop-projection-formula}
\dcref{B.16}\source{kleiman-picard:page-0076}
For an integral $Y'\subset X'$ mapping to $Y=g(Y')\subset X$,
\[
 \int c_1(g^*\mathcal L_1)\cdots c_1(g^*\mathcal L_r)[Y']
 =\deg(Y'/Y)\int c_1(\mathcal L_1)\cdots c_1(\mathcal L_r)[Y],
\]
where the degree is zero when the function-field extension is not finite.
\uses{kps-lem-intersection-pushforward,kps-lem-cycle-decomposition,kps-thm-first-intersection}
\end{proposition}
\begin{proof}
The cycle decomposition of $g_*\mathcal O_{Y'}$ is
$\deg(Y'/Y)[Y]$ modulo lower-dimensional classes. Apply Lemma B.15 and
Theorem B.9(1).
\end{proof}

\begin{proposition}[Invariance under field extension]\label{kps-prop-intersection-base-change}
\dcref{B.17}\source{kleiman-picard:page-0076}
Intersection numbers are unchanged after extending the ground field.
\uses{kps-def-K,kps-def-intersection-number}
\end{proposition}
\begin{proof}
Flat base change preserves short exact sequences and Euler characteristics,
therefore induces a map on filtered Grothendieck groups preserving all the
$c_1$ operators and $\chi$.
\end{proof}

\begin{proposition}[Constancy in flat families]\label{kps-prop-intersection-family}
\dcref{B.18}\source{kleiman-picard:page-0077}
If $\mathcal F$ is flat with proper support of relative dimension $r$ over
$S$, then
$y\mapsto\int c_1(\mathcal L_1)\cdots c_1(\mathcal L_r)\mathcal F_y$
is locally constant.
\uses{kps-def-intersection-number}
\end{proposition}
\begin{proof}
The intersection number is an Euler characteristic of a flat proper family;
its local constancy is the standard constancy theorem for such Euler
characteristics.
\end{proof}

\begin{definition}[Numerical equivalence in families]\label{kps-def-family-numerical}
\dcref{B.19}\source{kleiman-picard:page-0077}
Line bundles $\mathcal L,\mathcal N$ are numerically equivalent when
$\int c_1(\mathcal L)[Y]=\int c_1(\mathcal N)[Y]$ for every closed integral
curve $Y$ with proper support over an Artin subscheme of $S$.
\uses{kps-def-intersection-number}
\end{definition}

\begin{proposition}[Numerical invariance of intersections]\label{kps-prop-numerical-invariance}
\dcref{B.20}\source{kleiman-picard:page-0077}
If $\mathcal L_j$ and $\mathcal N_j$ are numerically equivalent for every $j$,
then their $r$-fold intersections against any $\mathcal F\in\mathcal K_r$
are equal.
\uses{kps-def-family-numerical,kps-lem-dimension-drop,kps-lem-intersection-cycles}
\end{proposition}
\begin{proof}
For $r=1$ use the intersection-cycle lemma. For larger $r$, dimension drop
puts the remaining product in $\mathcal K_1$; replace the factors one at a
time using numerical equivalence and symmetry.
\end{proof}

\begin{proposition}[Numerical equivalence and pullback]\label{kps-prop-numerical-pullback}
\dcref{B.21}\source{kleiman-picard:page-0077}
Pullback preserves numerical equivalence. If $g:X'\to X$ is proper and
surjective, numerical equivalence of pullbacks implies numerical equivalence
on $X$.
\uses{kps-prop-projection-formula}
\end{proposition}
\begin{proof}
Apply the projection formula to every integral curve upstairs. Conversely,
lift a closed integral curve downstairs to a curve upstairs mapping onto it;
the resulting nonzero degree lets the projection-formula equalities descend.
\end{proof}

\begin{definition}[Degree on proper curves]\label{kps-def-degree}
\dcref{B.22}\source{kleiman-picard:page-0077}
For an Artin base and a proper curve, define
$\deg(\mathcal L)=\int c_1(\mathcal L)$ and
$\deg(D)=\deg(\mathcal O_X(D))$.
\uses{kps-def-intersection-number}
\end{definition}

\begin{proposition}[Riemann's theorem for curves]\label{kps-prop-riemann}
\dcref{B.23}\source{kleiman-picard:page-0078}
Degree is additive; an effective divisor satisfies
$\deg(D)=h^0(\mathcal O_D)$; and
\[
 \chi(\mathcal L)=\deg(\mathcal L)+\chi(\mathcal O_X).
\]
For an integral curve, degree is unchanged by normalization pullback.
\uses{kps-thm-first-intersection,kps-lem-intersection-divisor,kps-prop-projection-formula}
\end{proposition}
\begin{proof}
Additivity is Theorem B.9(2), the divisor formula is Lemma B.13, and the
Riemann formula follows from
$\deg(\mathcal L^{-1})=\chi(\mathcal O_X)-\chi(\mathcal L)$. Normalization
invariance is the projection formula.
\end{proof}

\begin{definition}[Arithmetic genus of a divisor]\label{kps-def-arithmetic-genus}
\dcref{B.24}\source{kleiman-picard:page-0078}
For a divisor $D$ on a proper surface over an Artin base, set
\[
 p_a(D)=1-\chi(c_1(\mathcal O_X(D))\mathcal O_X).
\]
\uses{kps-def-intersection-number}
\end{definition}

\begin{proposition}[Genus and intersection formula]\label{kps-prop-arithmetic-genus}
\dcref{B.25}\source{kleiman-picard:page-0078}
For divisors $D,E$ on a proper surface,
$p_a(D+E)=p_a(D)+p_a(E)+(D\cdot E)-1$. If $D$ is effective, then
$p_a(D)=1-\chi(\mathcal O_D)$.
\uses{kps-lem-c1-additivity,kps-lem-divisor-action}
\end{proposition}
\begin{proof}
Expand $c_1(\mathcal O(D+E))$ using Lemma B.3 and apply the divisor-action
lemma to the effective case.
\end{proof}

\begin{proposition}[Riemann--Roch for surfaces]\label{kps-prop-surface-RR}
\dcref{B.26}\source{kleiman-picard:page-0078}
Let $X$ be a reduced projective equidimensional Cohen--Macaulay surface over
a field, with dualizing sheaf $\omega$ and $\mathcal K=\omega-\mathcal O_X$.
Then $\mathcal K\in\mathcal K_1$ and, for every divisor $D$,
\[
 p_a(D)=\frac{(D^2)+(D\cdot\mathcal K)}2+1,\qquad
 \chi(\mathcal O_X(D))=\frac{(D^2)-(D\cdot\mathcal K)}2+\chi(\mathcal O_X).
\]
If $X$ is Gorenstein with canonical divisor $K$, replace $\mathcal K$ by
$K$ in the intersection expressions.
\uses{kps-thm-first-intersection}
\end{proposition}
\begin{proof}
The dualizing sheaf agrees with $\mathcal O_X$ on a dense open, so its
difference lies in $\mathcal K_1$. Expand the two sides using Theorem B.9 and
use duality $\chi(\mathcal L)=\chi(\mathcal L^{-1}\ox\omega)$; solving the
resulting two linear equations gives the displayed formulas. In the
Gorenstein case $\mathcal K=-c_1(\omega^{-1})\mathcal O_X$, yielding $K$.
\end{proof}

\begin{theorem}[Hodge index theorem]\label{kps-thm-hodge-index}
\dcref{B.27}\source{kleiman-picard:page-0079}
Let $X$ be a geometrically irreducible complete surface over a field and
$\mathcal H,\mathcal L\in\Pic(X)$. If
$\int c_1(\mathcal H)^2>0$,
$\int c_1(\mathcal L)c_1(\mathcal H)=0$, and
$\int c_1(\mathcal L)^2\ge0$, then $\mathcal L$ is numerically equivalent to
$\mathcal O_X$.
\end{theorem}
\begin{proof}
Extend the field, reduce, and pass by Chow's lemma and normalization to an
integral projective surface using Propositions B.17, B.16, and B.21. If a
curve $Y$ had nonzero $\mathcal L$-degree, blow it up and use the exceptional
divisor to obtain a line bundle $\mathcal M$ with
$\int c_1(\mathcal L)c_1(\mathcal M)\ne0$. Ample twists and powers then give
bundles $\mathcal L_1,\mathcal H_1$ with $\mathcal H_1$ ample,
$\mathcal L_1\mathcal H_1=0$, and $\mathcal L_1^2>0$. Choose
$\mathcal N=\mathcal L_1^n\ox\mathcal H_1^{-1}$ with positive intersection
against an ample twist. Positivity (Proposition B.14) forces simultaneously
$\mathcal N\mathcal H_1<0$ and $\mathcal N^2>0$, contradicting Lemma B.28.
\uses{kps-prop-intersection-base-change,kps-prop-projection-formula,kps-prop-numerical-pullback,kps-prop-ample-positivity,kps-lem-intersection-divisor,kps-lem-positive-square}
\end{proof}

\begin{lemma}[Positive-square criterion]\label{kps-lem-positive-square}
\dcref{B.28}\source{kleiman-picard:page-0080}
For an integral surface and $\mathcal N$ with $\int c_1(\mathcal N)^2>0$,
the following are equivalent: $\mathcal N$ has positive intersection with
every ample bundle; it has positive intersection with some ample bundle; and
$H^0(\mathcal N^n)\ne0$ for some $n>0$.
\uses{kps-prop-ample-positivity,kps-thm-snapper}
\end{lemma}
\begin{proof}
An effective divisor in a positive power of $\mathcal N$ intersects every
ample bundle positively by Proposition B.14. Conversely, if an ample bundle
has positive intersection with $\mathcal N$, duality and the Snapper
polynomial show $H^2(\mathcal N^n)=0$ for large $n$, while the positive
quadratic term in $\chi(\mathcal N^n)$ forces $H^0(\mathcal N^n)\ne0$.
\end{proof}

\begin{corollary}[Connectedness of ample divisors]\label{kps-cor-connected-divisors}
\dcref{B.29}\source{kleiman-picard:page-0080}
If $X$ is a geometrically irreducible projective $r$-fold with $r\ge2$, $D$
is ample, and $E$ is effective, then $D+E$ is connected.
\end{corollary}
\begin{proof}
After replacing $D$ by a suitable ample $nD+E$, suppose it split as disjoint
$D_1+D_2$. For $r=2$, disjointness gives $(D_1\cdot D_2)=0$, while positivity
gives $(D_i^2)>0$; this contradicts the Hodge index theorem. For $r\ge3$, a
general integral hyperplane section preserves the disjoint decomposition and
reduces to the lower-dimensional case.
\uses{kps-lem-intersection-divisor,kps-prop-ample-positivity,kps-thm-hodge-index}
\end{proof}
